%!TEX program = xelatex
% 完整编译: xelatex -> biber/bibtex -> xelatex -> xelatex
\documentclass[lang=cn,a4paper]{shu-lab-report}

\usepackage{standalone}
% 本文档命令
\usepackage{array}
\newcommand{\ccr}[1]{\makecell{{\color{#1}\rule{1cm}{1cm}}}}
\renewcommand{\lstlistingname}{代码}




\begin{document}

%%%%%%%%%%%%%%%%%%%%%%%%%%%%%
%% The Cover Page
%%%%%%%%%%%%%%%%%%%%%%%%%%%%%

\thispagestyle{empty}

~ \\

\vspace{3cm}

\begin{figure}[!htbp]
    \centering
    
\includegraphics[width=10cm]{image/shulogo.png}
\end{figure}

\centerline{\large{\textbf{SHANGHAI  UNIVERSITY}}}

\vspace{8mm}

\centerline{\kaishu\Huge{\textbf{机器学习基础~实验报告}}}

\vspace{8mm}


\vspace{16mm}

\begin{center}
\renewcommand\arraystretch{1.5}
\begin{tabular}{r c}
    \makebox[8em][s]{\LARGE{学号姓名}} & \LARGE{21122915王佳信}\\  \cmidrule(l){2-2} 
    \makebox[8em][s]{\LARGE{实验序号}} & \LARGE{实验4~贝叶斯分类器}\\            \cmidrule(l){2-2} 
    \makebox[8em][s]{\LARGE{日期}}    & \LARGE{2023年12月18日}\\ \cmidrule(l){2-2}
\end{tabular}
\end{center}

\vspace{8mm}

%%%%%%%%%%%%%%%%%%%%%%%%%%%%%
%% The Content
%%%%%%%%%%%%%%%%%%%%%%%%%%%%%

\newpage
\setcounter{page}{1}

\section{实验目的与要求}
目的：熟悉和构建贝叶斯分类

内容：
\begin{itemize}
    \item 构建贝叶斯分类器，并计算准确率
    \item 对比贝叶斯分类器和决策树，及线性模型
\end{itemize}

\section{实验环境}
\textbf{工具}: python, numpy, pandas, matplotlib

\textbf{数据集}: 
男生女生身高体重鞋码数据集。

数据集中包括男生女生身高体重鞋码数据，数据集中共有800个样本，每个样本包括4个属性，分别是身高、体重、鞋码和性别。其中身高和体重是连续属性，鞋码是离散属性，性别是类别属性。


\section{实验内容及其设计与实现}

\subsection{实验内容1~贝叶斯分类器}
\subsubsection{算法原理}
朴素贝叶斯是经典的机器学习算法之一，也是为数不多的基于概率论的分类算法。
对于大多数的分类算法，在所有的机器学习分类算法中，
朴素贝叶斯和其他绝大多数的分类算法都不同。比如决策树，KNN，逻辑回归，支持向量机等，
他们都是判别方法，也就是直接学习出特征输出Y和特征X之间的关系，要么是决策函数，
要么是条件分布。但是朴素贝叶斯却是生成方法，该算法原理简单，也易于实现。

朴素贝叶斯算法的核心思想是基于贝叶斯定理和特征条件独立假设，即假设特征之间相互独立。
贝叶斯公式便是基于条件概率，通过 $P(B|A)$ 来求 $P(A|B)$ :
\begin{equation}
    P(A|B) = \frac{P(B|A)P(A)}{P(B)}
\end{equation}

朴素贝叶斯对条件概率分布作出了独立性的假设，在这个假设的前提上，条件概率可以转化为：
\begin{equation}
    P(x|y_k)=P(x_1,x_2,...,x_n|y_k)=\prod_{i=1}^{n}{P(x_i|y_k)}
\end{equation}

代入上面贝叶斯公式中，得到：
\begin{equation}
    P(y_k|x)=\frac{P(y_k)\ast\prod_{i=1}^{n}{P(x_i|y_k)}}{\sum_{k}{P(y_k)}\ast\prod_{i=1}^{n}{P(x_i|y_k)}}
\end{equation}

于是，朴素贝叶斯分类器可表示为：
\begin{equation}
    y=f(x)=\arg\max_{y_k}{P(y_k)\ast\prod_{i=1}^{n}{P(x_i|y_k)}}
\end{equation}



朴素贝叶斯算法的基本步骤：
\begin{enumerate}
    \item 计算先验概率及类条件概率密度
    \item 对于给定的样本特征，通过贝叶斯公式计算后验概率
    \item 若按最小错误率，选择具有最大后验概率的类别作为预测结果，
    或最小风险原则分类
\end{enumerate}

\subsubsection{对比贝叶斯分类器和决策树、线性模型}
\textbf{贝叶斯分类器}：基于概率论的分类器，对于连续特征，其假设特征服从正态分布，
计算均值和标准差，代入概率密度函数，计算后验概率，进行分类。
对于离散特征，其假设特征服从多项式分布，计算每个特征的概率，代入后验概率公式，进行分类。

朴素贝叶斯算法的优点：
\begin{itemize}
    \item 朴素贝叶斯模型有稳定的分类效率。
    \item 对小规模的数据表现很好，能处理多分类任务，适合增量式训练，
    尤其是数据量超出内存时，可以一批批的去增量训练。
    \item 对缺失数据不太敏感，算法也比较简单，常用于文本分类。
\end{itemize}

朴素贝叶斯算法的缺点：
\begin{itemize}
    \item 需要知道先验概率，且先验概率很多时候取决于假设，
    假设的模型可以有很多种，因此在某些时候会由于假设的先验模型的原因导致预测效果不佳。
    \item 对输入数据的表达形式很敏感（0值、离散、连续，值极大极小之类的）。
\end{itemize}

\textbf{决策树}：基于信息论的分类器，十分适合离散特征的处理。
对于离散特征，其它很多模型都需要对其进行编码，得到很稀疏的编码向量再进行模型拟合，
因较大的过拟合风险对样本量有较高要求。决策树则统计离散特征后直接计算信息增益等，
再加上有各类集成树模型的存在,可以进一步提升性能。

决策树的优点是易于理解和解释，可以处理非线性和缺失数据，
可以生成可视化的树形结构。

决策树的缺点是容易过拟合，对噪声数据敏感，
可能产生不稳定的结果。决策树适用于特征维度较高，
数据量较大，类别分布较不平衡的情况。

\textbf{线性模型}：线性模型是一种基于函数的分类模型，
它通过构造一个线性函数来划分数据空间，使得不同类别的数据位于不同的区域。
线性模型的优点是计算效率高，可以进行特征选择和正则化，可以处理大规模的数据。
线性模型的缺点是不能处理非线性的数据，对异常值和多重共线性敏感，
可能存在欠拟合的风险。线性模型适用于特征维度较高，
数据量较大，数据分布较线性的情况。

\subsubsection{代码及结果}
对于单个特征的朴素贝叶斯分类器，假设每个特征都服从正态分布，根据正态分布的概率密度函数：
\begin{equation}
    f(x)=\frac{1}{\sigma\sqrt{2\pi}}e^{-\frac{(x-\mu)^2}{2\sigma^2}}
\end{equation}

计算出每个特征的均值和标准差，然后代入概率密度函数，计算出后验概率，进行分类。
其代码如下所示。
\begin{lstlisting}[language=Python, caption={朴素贝叶斯分类器}, label={code:naive_bayes}]
    def gaussian_pdf(x , mean , std):
    '''
    正态分布概率密度函数
    '''
    prob_density = (1/std*np.sqrt(2*np.pi)) * np.exp(-((x - mean) ** 2) / (2 * std ** 2))
    return prob_density

def bayes_classifier_by(train_data, test_data, feature='身高(cm)', target='性别', prior_prob_male=0.5):
    '''
    单特征特征的Bayes分类器
    '''
    prior_prob_female= 1 - prior_prob_male
    # 提取训练集中男女身高的均值和标准差
    male_mean = train_data[train_data[target] == '男'][feature].mean()
    male_std = train_data[train_data[target] == '男'][feature].std()
    female_mean = train_data[train_data[target] == '女'][feature].mean()
    female_std = train_data[train_data[target] == '女'][feature].std()
    # 计算男女身高的概率密度函数
    male_pdf = gaussian_pdf(test_data[feature], male_mean, male_std)
    female_pdf = gaussian_pdf(test_data[feature], female_mean, female_std)
    # 计算后验概率
    posterior_prob_male = prior_prob_male * male_pdf
    posterior_prob_female = prior_prob_female * female_pdf
    # 进行分类
    predictions = np.where(posterior_prob_male > posterior_prob_female, '男', '女')
    return predictions

\end{lstlisting}

结果性能分析（以ROC曲线作为分析指标）如下：
\begin{figure}[!htbp]
    \centering
    \includegraphics[width=0.32\textwidth]{image/single_feature_output.png}
    \caption{不同特征ROC曲线}
    \label{fig:roc}
\end{figure}

同时采用身高和体重数据作为特征，采用最小错误率原则进行分类，其代码如下所示。
\begin{lstlisting}[language=Python, caption={最小错误率分类器}, label={code:minimum_error_rate}]
    class MyGaussianNB:
    def fit(self, X, y, prior_probs=None):
        self.classes, class_counts = np.unique(y, return_counts=True)
        if prior_probs is not None:
            prior_probs = {'男': prior_probs, '女': 1 - prior_probs}
        else:
            prior_probs = dict(zip(self.classes, class_counts / len(y)))
        self.class_probs = np.array([prior_probs[class_label] for class_label in self.classes])
        self.means = {}
        self.stds = {}
        for class_label in self.classes:
            class_data = X[y == class_label]
            self.means[class_label] = np.mean(class_data, axis=0)
            self.stds[class_label] = np.std(class_data, axis=0)

    def predict(self, X):
        predictions = []
        for x in X:
            class_scores = []
            for i, class_label in enumerate(self.classes):
                prior_prob = np.log(self.class_probs[i])
                mean = self.means[class_label]
                std = self.stds[class_label]
                likelihood = np.sum(-0.5 * ((x - mean) / std) ** 2 - np.log(np.sqrt(2 * np.pi) * std))
                class_scores.append(prior_prob + likelihood)
            predictions.append(self.classes[np.argmax(class_scores)])   #最大错误率
        return np.array(predictions)
    
    def predict_proba(self, X):
        probas = []
        for x in X:
            class_probs = []
            for i in range(len(self.classes)):
                class_label = self.classes[i]
                prior_prob = np.log(self.class_probs[i])
                mean = self.means[class_label]
                std = self.stds[class_label]
                likelihood = np.sum(-0.5 * ((x - mean) / std) ** 2 - np.log(np.sqrt(2 * np.pi) * std))
                class_probs.append(np.exp(prior_prob + likelihood))
            class_probs /= np.sum(class_probs)  # Normalize to get probabilities
            probas.append(class_probs)
        return np.array(probas)
\end{lstlisting}

结果性能分析（以ROC曲线作为分析指标）如下：
\begin{figure}[!htbp]
    \centering
    \includegraphics[width=0.4\textwidth]{image/multifeature.png}
    \caption{多特征ROC曲线}
    \label{fig:roc}
\end{figure}

三维高斯分布及决策面投影可视化：
\begin{figure}[!htbp]
    \centering
    \begin{minipage}[t]{0.48\textwidth}
        \centering
        \includegraphics[width=6cm]{image/3d.png}
        \caption{三维高斯分布}
    \end{minipage}
    \begin{minipage}[t]{0.48\textwidth}
        \centering
        \includegraphics[width=6cm]{image/decision_boundary_output.png}
        \caption{决策面投影}
    \end{minipage}
    \label{fig:visualizaion}
\end{figure}

\section{实验内容2~Parzen窗法}
\subsection{算法原理}
Parzen窗法是一种非参数密度估计方法，其基本思想是将一个窗函数放在每个样本点上，
然后通过计算窗函数在各个样本点上的积分来估计概率密度函数。
具体来说，Parzen窗法将样本点分为若干个矩形区域，然后计算每个区域内样本点的个数，
最后通过这些个数来估计概率密度函数。
公式如下：
\begin{equation}
    \hat{p}(x)=\frac{1}{n}\sum_{i=1}^{n}\frac{1}{h^d}K(\frac{x-x_i}{h})
\end{equation}

核函数为：
\begin{equation}
    \text{{kernel\_sum}} = \frac{1}{n} \sum_{i=1}^{n} \frac{e^{-\frac{(x - x_i)^2}{2h^2}}}{\sqrt{2\pi}h}
\end{equation}

\subsection{代码与结果}
代码如下所示：
\begin{lstlisting}[language=Python, caption={Parzen窗法}, label={code:parzen}]
def parzen_window_pdf(x, data, h):
    '''
    Parzen窗法概率密度估计函数
    '''
    n = len(data)
    prob_density = np.sum(np.exp(-(x - data) ** 2 / (2 * h ** 2)) / (np.sqrt(2 * np.pi) * h)) / n
    return prob_density

def bayes_classifier_parzen_height(train_data, test_data, h=5.0, prior_prob_male=0.5, prior_prob_female=0.5):
    '''
    以身高为特征的Bayes分类器（Parzen窗法）
    '''
    predictions = []
    for index, test_point in test_data.iterrows():
        # 提取训练集中男女身高数据
        male_height_data = train_data[train_data['性别'] == '男']['身高(cm)']
        female_height_data = train_data[train_data['性别'] == '女']['身高(cm)']
        # 计算男女身高的概率密度函数
        male_height_pdf = parzen_window_pdf(test_point['身高(cm)'], male_height_data, h)
        female_height_pdf = parzen_window_pdf(test_point['身高(cm)'], female_height_data, h)
        # 计算后验概率
        posterior_prob_male = prior_prob_male * male_height_pdf
        posterior_prob_female = prior_prob_female * female_height_pdf
        # 进行分类
        prediction = '男' if posterior_prob_male > posterior_prob_female else '女'
        predictions.append(prediction)
    return predictions
\end{lstlisting}

结果性能分析（以不同窗宽下ROC曲线作为分析指标）如下：
\begin{figure}[!htbp]
    \centering
    \includegraphics[width=0.4\textwidth]{image/Parzen.png}
    \caption{Parzen窗法ROC曲线}
    \label{fig:roc}
\end{figure}


\section{总结与体会}
通过本次实验，我对贝叶斯分类器有更深入的理解，对其公式推导有了更深刻的认识，
同时也对Parzen窗法的原理有了更深入的了解。

\end{document}




\section{模板测试部分（实际编写时请删除）}

如\coderef{code:example}所示，为示例代码。
\begin{lstlisting}[language=Python, caption={示例代码}, label={code:example}]

\end{lstlisting}


如\figref{fig:example}所示，为上海大学校徽。
\begin{figure}[!htbp]
    \centering
    
\includegraphics[width=0.5\textwidth]{image/shulogo.png}
    \caption{上海大学校徽}
    \label{fig:example}
\end{figure}

如\tabref{tab:example}所示，为测试表。
\begin{table}[!htbp]
    \caption{测试表}
    \label{tab:example}
    \centering
    \begin{tabular}{cc}
        \toprule
            & value\\
        \midrule
        1   & 1     \\
        2   & 2     \\
        \bottomrule
    \end{tabular}
\end{table}


\textbf{说明}：叙述各个实验任务的过程和结果，着重说明各个实验所反映的
操作系统设计原理，加以评述。

\begin{itemize}
    \item 
    \item 
    \item 
\end{itemize}

\begin{enumerate}
    \item 
    \item 
    \item 
\end{enumerate}